\section{A lower bound for the continuum energy}\label{sec:dual}

For a measure $\sigma$ write
$L_\sigma(z)=\int\log|z-w|\,d\sigma(w)$. Fix an admissible reference
pair $(\eta_0,\mu_0)$ with $\eta_0=\kappa_0+\nu_0$. Its exact cell masses
and contour weights are given in the supplementary certificate files.
Define
\[
\begin{aligned}
Q_{\log}={}&\iint\log|z-w|\,d\eta_0(z)d\eta_0(w)
 +\iint\log|t-s|\,d\mu_0(t)d\mu_0(s)\\
 &+\iint\log|z+t|\,d\eta_0(z)d\mu_0(t),\\
W_\eta={}&A_0-\phi-2L_{\eta_0}-L_{\mu_0^-},\\
W_\mu(t)={}&B_0(t)-\psi(t)-2L_{\mu_0}(t)-L_{\eta_0}(-t).
\end{aligned}
\]
For constants $h_\eta,h_\mu$ satisfying $W_\eta\ge h_\eta$ on $\Gamma$,
put
\[
P=\int_E(h_\eta-W_\eta)_+\,dx,
\qquad T=\int_{1/2}^{\infty}(h_\mu-W_\mu)_+\,dt.
\]

\begin{lemma}[Dual lower bound]\label{lem:dual-lower}
For every admissible pair of measures,
\begin{equation}
\alpha_\Gamma\ge C_0+Q_{\log}+h_\eta+h_\mu-P-T.
\label{eq:duality-bound}
\end{equation}
\end{lemma}

\begin{proof}
Let $\sigma=\eta-\eta_0$ and $\tau=\mu^- -\mu_0^-$; both have total
mass zero. The frequency variation is written on the reflected negative
axis. Reflection preserves its self-energy, so $I(\tau)=I(\mu-\mu_0)$.
The quadratic remainder after linearization is
\[
I(\sigma)+I(\tau)+I(\sigma,\tau)
 =\tfrac12 I(\sigma)+\tfrac12 I(\tau)+\tfrac12 I(\sigma+\tau)\ge0.
\]
A signed measure of mass zero has nonnegative logarithmic energy. For
smooth measures this is the Dirichlet integral identity. Finite-energy
measures follow by approximation. In the frequency first variation, the
change of variables from $x<0$ to $t=-x$ gives $W_\mu(t)$. Expanding the
energy and using the definitions of $W_\eta,W_\mu$ gives
\[
\mathcal E\ge C_0+Q_{\log}+\int W_\eta\,d\eta+\int W_\mu\,d\mu.
\]
The pinned and contour masses sum to one, so the first integral is at
least $h_\eta-P$. By the frequency cap the second integral is at least
$h_\mu-T$. This proves \eqref{eq:duality-bound}.
\end{proof}

The reference has pinned and frequency densities in $[0,1]$ and the
required total masses. On the contour it is a nonnegative symmetric
Fej\'er mixture. Its cell masses and Fej\'er weights are rational. The
checker evaluates the finitely many algebraic Fourier moments by directed
interval enclosures. The supplementary materials give the complete
reference, the potential formulas and closed endpoint evaluations. The
comparison uses no quadrature value.

Take $h_\eta=-669/200$ and $h_\mu=71/25$. The directed-interval
certificate in the supplementary materials checks the bounds on the full
contour, pinned and frequency covers. It also checks the exact reference
integrals:
\[
P\le\frac{6905963574153}{81920000000000},\quad
T_{[1/2,81/10]}\le\frac{38158094390353}{409600000000000},\quad
Q_{\log}>-\frac{438334707}{10^9}.
\]
The rational cell lists and the directed checker are in the supplementary
archive. For the infinite frequency tail we argue by hand. For
$t\ge81/10$,
\[
W_\mu'(t)\ge\log2-\frac1{2(t-1/2)}+\frac3{t+2}
 -\frac2{t-11/2}-\frac1{t-751/5000}.
\]
The right side is increasing and exceeds $1/50$ at $81/10$. The checker
also verifies
\[
W_\mu(81/10)>71/25=h_\mu.
\]
Thus $W_\mu$ is increasing above this endpoint, and
$(h_\mu-W_\mu(t))_+=0$ for every $t\ge81/10$. So the positive-part
integral $T$ has no tail beyond that point. A rational upper bound for
$\log2$ in \eqref{eq:duality-bound} now gives
\begin{equation}
\alpha_\Gamma>\frac{1992909}{10^6}.
\label{eq:alpha-lower}
\end{equation}
